\documentclass[12pt,a4paper]{article}

%pacotes necessários
\usepackage[utf8]{inputenc}
\usepackage[T1]{fontenc}
\usepackage[brazilian]{babel}
\usepackage{geometry}
\usepackage{graphicx}
\usepackage{amsmath}
\usepackage{listings}
\usepackage{xcolor}
\usepackage{tikz}
\usetikzlibrary{automata, positioning, arrows.meta}
\usepackage{float}
\usepackage{booktabs}
\usepackage[hidelinks]{hyperref}
\input{glyphtounicode}
\pdfgentounicode=1

%margens 
\geometry{
    left= 3cm,
    right= 2cm,
    top= 3cm,
    bottom= 2cm
}

\begin{document}
\hypersetup{pageanchor=false}
%capa
\begin{titlepage}
\centering

{\large CENTRO UNIVERSITÁRIO MAURÍCIO DE NASSAU}\\[1cm]

{\large CURSO DE CIÊNCIA DA COMPUTAÇÃO}\\[3cm]

{\Large\textbf{ANALISADOR LÉXICO (SCANNER)}}\\[2cm]

Alessandra Júlia Santos Mendonça - 01710330\\
Jhonas Xavier dos Santos SILVA - 01694333\\
Marcos Emanuel Celestino Tavares - 01718690\\
Guilherme de Oliveira Lopes Silva - 01699872\\
Alan Sergio Rodrigues de Carvalho - 01697157\\
Alfred Manoel Vincente Silva da Rocha - 01695802\\
Pedro Artur Rolemberg Ferreira Nogueira - 01705061\\

\vfill
{\large Maceió - AL}\\
{\large 2026}

\end{titlepage}
\hypersetup{pageanchor=true}
\pagenumbering{roman}
%sumário
\tableofcontents
\clearpage
\pagenumbering{arabic}


\section{Introdução}

O desenvolvimento de uma mini-linguagem de programação envolve a definição de uma estrutura própria de representação de comandos, estabelecendo regras que determinam como as instruções devem ser escritas e interpretadas. Nesse contexto, o presente projeto propõe a criação de uma mini-linguagem fundamentada no paradigma de programação orientada a eventos, cuja sintaxe e semântica são inspiradas em comportamentos e ações associados aos ratos.

A proposta consiste em estabelecer uma linguagem temática na qual os elementos utilizados na construção dos comandos estejam relacionados a situações e comportamentos desses animais. Dessa maneira, conceitos como eventos e ações servem como referência para a organização das instruções, proporcionando uma identidade própria à linguagem e uma abordagem diferenciada para a aplicação dos conceitos de linguagens formais.

No âmbito técnico, o relatório apresenta os conceitos e as regras do analisador léxico (Scanner), responsável por percorrer o código-fonte e classificar seus elementos em categorias denominadas tokens. A especificação formal ainda registra decisões pendentes de validação pelo grupo.

Na execução de arquivos \texttt{.ra}, o programa usa Scanner, Parser e Executor. O console interativo segue um caminho separado e traduz algumas palavras para Python. O relatório aborda linguagens formais, expressões regulares e autômatos usados no Scanner.

\section{Objetivos}

Desenvolver e documentar a especificação de um analisador léxico (Scanner) para uma linguagem de programação temática, utilizando linguagens formais, expressões regulares e autômatos finitos determinísticos para reconhecer e classificar os elementos léxicos.

\subsection{Objetivo Geral}
\begin{itemize}

    \item Definir o alfabeto e as categorias de tokens da linguagem;
    \item Estabelecer expressões regulares para representar os padrões léxicos;
    \item Modelar autômatos finitos determinísticos para o reconhecimento de identificadores, números e operadores;
    \item Descrever o processo de identificação e classificação dos lexemas; 
    \item Aplicar a estratégia de maior casamento (\textit{maximal munch}) durante o reconhecimento dos tokens;
    \item Especificar mecanismos para identificação e tratamento de erros léxicos;
    \item Documentar as simulações da proposta de Scanner.

\end{itemize}

    
\subsection{Objetivos Específicos}
\begin{itemize}
\item Definir as categorias de tokens da linguagem.
\item Representar os padrões com expressões regulares e autômatos.
\item Descrever o tratamento de erros léxicos.
\end{itemize}

\section{Fundamentação Teórica}
\subsection{Compiladores e Análise Léxica}
Um compilador é um programa responsável por traduzir um código escrito em uma linguagem de programação para outra representação, geralmente uma linguagem de máquina ou uma forma intermediária. Esse processo é dividido em etapas, que analisam diferentes aspectos do código-fonte.

A análise léxica é uma das primeiras etapas desse processo. Sua responsabilidade é percorrer o código fonte, identificar sequências de caracteres que formam unidades significativas e classificá-las em categorias denominadas tokens.

O componente responsável por essa atividade é chamado de analisador léxico ou scanner. Além de reconhecer os elementos válidos da linguagem, ele também pode identificar caracteres inválidos e informar a ocorrência de erros léxicos.

\subsection{Tokens e Lexemas}

Token é a classificação atribuída a uma unidade reconhecida durante a análise léxica. Já o lexema corresponde à sequência de caracteres presentes no código-fonte que representa determinado token. 

Por exemplo, considerando a instrução

\begin{lstlisting}
    roer queijo = 18;
\end{lstlisting}

Na proposta, a instrução abaixo seria dividida nos lexemas \texttt{roer}, \texttt{queijo}, \texttt{=}, \texttt{18} e \texttt{;}, cada um associado à sua categoria léxica.

\subsection{Expressões Regulares}

Expressões regulares são representações formais utilizadas para descrever padrões de cadeias de caracteres. Elas permitem definir quais sequências pertencem a determinadas linguagens regulares.

No desenvolvimento do Scanner, as expressões regulares são utilizadas para especificar os padrões de reconhecimento dos tokens, como identificadores, números inteiros, números decimais e operadores. 

\subsection{Autômatos Finitos}

Autômatos finitos são modelos matemáticos utilizados para reconhecer linguagens regulares. Um autômato finito determinístico é definido por um conjunto de estados, um alfabeto, uma função de transição, um estado inicial e um conjunto de estados finais. 

\section{Metodologia}
Foram reunidas as palavras e os símbolos previstos para a linguagem. Depois, foram definidos padrões para os tokens e implementados reconhecedores por meio de autômatos. Os testes do projeto verificam o Scanner, o Parser e casos de integração.
\section{Desenvolvimento do Scanner}
\subsection{Especificação da Linguagem e Categorias de Tokens}

A proposta considera palavras reservadas, identificadores, números inteiros e decimais, strings entre aspas duplas, operadores e delimitadores. Colchetes são delimitadores usados em listas.

A linguagem também permite comentários de linha iniciados por \texttt{//}. O Scanner ignora todo o conteúdo após esse marcador até o final da linha, sem consumir a quebra de linha.

A lista completa e as decisões pendentes estão em \href{../markdown/SPECIFICATIONS.md}{SPECIFICATIONS.md}.

\subsection{Modelagem Matemática e Expressões Regulares}
Os padrões e as tabelas dos autômatos estão descritos em \href{../markdown/REGULAR_EXPRESSIONS.md}{expressões regulares} e \href{../markdown/AUTOMATA.md}{autômatos}.

Um identificador segue o padrão \texttt{(L|\_)(L|D|\_)*}, em que L representa uma letra e D, um dígito.

\subsection{Implementação e Funcionamento do Scanner}

O Scanner implementado em \texttt{src/lexer.py} reconhece tokens e registra suas posições no código. Strings são reconhecidas por uma regra própria, com aspas duplas e escapes definidos na especificação. Ele também reconhece as diretivas \texttt{\#tunel}, \texttt{\#toca} e \texttt{\#vaparatunel} por uma regra específica. O Parser, em \texttt{src/parser.py}, usa esses tokens para montar a árvore do programa. O Executor, em \texttt{src/executor.py}, percorre essa árvore e executa as instruções.

\subsection{Tratamento de Erros Léxicos}

Quando encontra caracteres fora do alfabeto definido, o Scanner registra erros com linha e coluna e continua a leitura. Ao analisar um arquivo, o Parser interrompe a execução se houver erro léxico ou sintático.

\section{Diagramas dos Autômatos}
As tabelas de transição e simulações dos modelos propostos estão em \href{../markdown/AUTOMATA.md}{AUTOMATA.md}.
\subsection{AFD de Identificadores}
O modelo e as simulações estão em \href{../markdown/AUTOMATA.md}{AUTOMATA.md}.
\subsection{AFD de Números}
O modelo e as simulações estão em \href{../markdown/AUTOMATA.md}{AUTOMATA.md}.
\subsection{AFD de Operadores Relacionais e Atribuição}
Os modelos e as simulações estão em \href{../markdown/AUTOMATA.md}{AUTOMATA.md}.
\subsection{AFDs de Operadores Aritméticos e Delimitadores}
Os modelos e as simulações estão em \href{../markdown/AUTOMATA.md}{AUTOMATA.md}.

\section{Testes e Resultados}
O projeto contém testes automatizados para o Scanner e o Parser, além de testes de integração. Eles verificam o reconhecimento de tokens, a formação da árvore e a execução do programa.
\section{Conclusão}
O trabalho apresenta as regras léxicas e os autômatos do Rathyon. No programa atual, o Scanner, o Parser e o Executor formam o caminho usado para executar arquivos \texttt{.ra}. A especificação léxica continua em rascunho enquanto algumas decisões aguardam validação do grupo.
\section{Referências}
\begin{itemize}
\item Projeto Rathyon. Código-fonte: \texttt{src/lexer.py}, \texttt{src/parser.py}, \texttt{src/executor.py} e \texttt{src/automata.py}.
\item Projeto Rathyon. Testes automatizados: diretório \texttt{tests/}.
\item Projeto Rathyon. Especificação léxica: \href{../markdown/SPECIFICATIONS.md}{SPECIFICATIONS.md}.
\item Projeto Rathyon. Autômatos e expressões regulares: \href{../markdown/AUTOMATA.md}{autômatos} e \href{../markdown/REGULAR_EXPRESSIONS.md}{expressões regulares}.
\item Projeto Rathyon. Gramática: \href{../markdown/grammar.md}{grammar.md}.
\end{itemize}


\end{document}
